\section{Introduction}\label{sec:introduction}

\begin{figure}[ht]
  \centering
  \includegraphics[width=0.60\linewidth, keepaspectratio]{figures/example-small.png}
  \caption{An example Flatland railway environment.}
  \label{fig:environment}
\end{figure}

Railway systems face numerous operational challenges in planning and execution,
and have been the subject of research for decades, particularly since the middle of the twentieth century.
By 2002, it was proven that developing train timetables for a single track belongs to the class NP-Hard \cite{}.
As modern railway systems grow,
both in terms of their networks and their passenger demand,
the real problems they face become more complex
and require more advanced solutions.

Tasked with solving such complex issues
and inspired by advancements in digital technology,
the Swiss Federal Railway (SBB) founded \textit{Flatland},
a global online railway scheduling and rescheduling competition
to crowdsource novel approaches to modern-day problems.
An example Flatland environment is shown in Figure~\ref{fig:environment}.
My research began in 2023 as an attempt to use answer set programming (ASP)
as an alternative to the machine learning approaches.
This culminated in a master thesis that outlined the creation of a toolkit
now known as \textit{flaspland}
that provides an interface between \textit{Flatland} written in Python
and ASP written in clingo.
The work also established standard fact formats for representing railway environments
and put forth simple baseline encodings for safely navigating trains.

The research continues today in the form of Ph.D. work
under the supervision of Prof. Dr. Torsten Schaub,
and includes a graduate-level railway scheduling course,
where students explore solutions to open research topics in this area,
and theses written by highly-motivated students
with a particular affinity for
research at the intersection between
knowledge representation and public transportation.
In this paper, I will: 
provide \textbf{background context} into the history of railway research, modern approaches, and the advantages of answer set programming; 
outline my \textbf{research goals} for the next few years; 
highlight the \textbf{current status} of this research;
share some \textbf{preliminary results};
and bring attention to \textbf{milestones} that lie on the horizon.


%
%%% Local Variables:
%%% mode: LaTeX
%%% TeX-master: "paper"
%%% End:
